\section{Conclusion and perspectives}
\label{sec:conclusion2}

\textbf{Synthesis of the three readings.} The arithmetic cascade of the PT active primes $\{3, 5, 7\}$, fixed by the fixed point $\mustar = 15$, admits three canonical realisations of complementary natures.
\begin{description}[nosep,leftmargin=*]
\item[Algebro-geometric reading (\ref{sec:courbe}).] The cascade is the unique algebraic curve in the Kontsevich--Norbury class satisfying conditions (A) polynomial holonomy, (B) activation threshold, and (C) self-consistency. Its genus is $\genus(\Cper) = \mustar - 1 = 14$.
\item[Differential reading (\ref{sec:soliton}).] The Fisher--Bianchi metric $\gPT$ induced by the cascade satisfies the asymptotic Ricci soliton equation $\Ric + \Hess(f) = \lambda(\mu)\, \gPT$ with $\lambda(\mu) = -1/\mu^{4} + O(1/\mu^{5})$, coefficient $-1$ exact.
\item[Spectral reading (\ref{sec:spectrales}).] The axiom $\sper = 1/2$ manifests in the Riemann--von Mangoldt residue via the triple identity $1/8 = \sper^{2}/2 = \lambdaH = \DeltaMaslov$ (K2), and the logarithmic derivative of the residual factor admits the exact decomposition $-d/ds \log R(s) = \sum_p \log p\,[1/(p^{s} - 1) + A_p\,p^{-s}]$ verified to $10^{-25}$ in $\Re(s) > 1$ (A6).
\end{description}

\textbf{Differentiated epistemic statuses.} The results of this article fall into three categories that should be carefully distinguished.
\begin{enumerate}[nosep,leftmargin=*]
\item \emph{Theorems proved unconditionally}: the classification theorem and the genus formula (\ref{thm:courbe_classification_en}, \ref{thm:genre_en}); the Ricci quasi-soliton theorem (\ref{thm:soliton_PT_en}).
\item \emph{Proved identities}: K2 (\ref{thm:K2_en}) is proved structurally as the exact coincidence of three independent sub-identities; A6 (\ref{thm:A6_en}) is an exact analytic identity in $\Re(s) > 1$, verified numerically to all available precisions.
\item \emph{Conjectural program}: the RH-PT conjecture (\ref{conj:RH_PT_en}) is an open research program, under investigation~\cite{PT_NewMath_Consolidation}. Its eventual resolution would entail the Riemann Hypothesis.
\end{enumerate}

\textbf{Articulation with the companion article.} This article presupposes, for its full intelligibility, the reading of~\cite{PT_Article1_Foundations}, where the seven structural theorems $\Tzero$--$\Tsix$ and the three proved bridge axioms $\BAthree$--$\BAfive$ are established. The geometric and spectral readings developed here are non-trivial extensions: they do not modify the arithmetic core of PT, but reveal its cohomological, differential, and spectral facets. A reverse reading (from this article to Article~1) is also possible: the topological genus 14 and the soliton equation provide \emph{a posteriori} constraints on the cascade that reinforce the uniqueness of the active-prime choice $\{3, 5, 7\}$.

\textbf{Perspectives.} Three immediate research directions emerge from this article.
\begin{itemize}[nosep,leftmargin=*]
\item \emph{Exhaustive computation of the invariants $\omega_{g,n}^{\Cper}$.} The Eynard--Orantin topological recursion allows, in principle, the computation of all invariants $\omega_{g,n}^{\Cper}$ of the persistence curve. The first computations (\cite[\S 8]{PT_GeoFlow_PersistenceCurve}) confirm the structural substitution $2\pi^{2}/3 \to (1/2) \sum \gammap{p}$ relative to Mirzakhani's Airy curve. Extending this computation to $\omega_{2,1}^{\Cper}$ and beyond would test new cascade identities.
\item \emph{Exact identity of the scale bridge $\lambda \leftrightarrow B_0$.} The relation $\lambda(\mu) \sim B_0(\mu)^{-8/3}$ is verified numerically to $0{,}4\,\%$. The \emph{exact} identity (up to a coefficient) between the Ricci residue $\lambda$ and the leading coefficient $B_0$ of topological recursion would constitute a strong algebro-differential bridge.
\item \emph{Cohomological promotion of K2.} The derivation of identity K2 from a cohomological invariant of the involution $D : \qplus \leftrightarrow \qminus$ (item K3 of the program~\cite{PT_NewMath_Consolidation}) would provide a unified proof of the three sub-identities, and a rigorous framework for their extension to dimensions other than $\mathbb R^{4}$.
\end{itemize}

\textbf{Links to other PT applications.} The readings presented here are in active dialogue with several domains of PT phenomenology, in particular cosmology (the Fisher--Bianchi metric is the PT analogue of Bianchi~I, cf.~\cite{companion_M3,companion_M5}), high-energy physics (the coupling $\lambdaH$ identified in K2 is the LHC-measurable Higgs self-coupling, cf.~\cite{companion_M5}), and analytic number theory (the RH-PT conjecture relates the PT survival operator to the Riemann zeros). None of these applications is developed here; they constitute the \emph{downstream} horizon of the present article.

\textbf{Editorial status.} The pair of articles~\cite{PT_Article1_Foundations}--present, together with the foundational series M1--M5, constitutes, to the author's knowledge, the most consolidated input material for an external peer review of the Theory of Persistence. The theorems T0--T6, the bridge axioms BA0--BA5, the derivation of $\alpha_{\rm EM}$, the persistence curve, the Ricci quasi-soliton, and the spectral identities K2 and A6 are presented here with their precise epistemic status. The open questions (RH-PT program, exact identity $\lambda \leftrightarrow B_0$, cohomological promotion of K2) are explicitly flagged and constitute the immediate research horizon of the program.
